\subsection{Cyclotomic extensions}

Let \(p\) be the characteristic exponent of \(\mathbf{K}\) and let \(n \geqslant 1\) be an integer prime to \(p\) ; by a cyclotomic extension of level \(n\) over \(\mathbf{K}\) we understand any splitting extension \(\mathbf{E}\) of the polynomial \(\mathbf{X}^n - 1\) over \(\mathbf{K}\) (V, p.~21). Since this polynomial is separable, \(\mathbf{E}\) is a Galois extension of \(\mathbf{K}\) , of finite degree (V, p.~57). There exists a primitive \(n\) -th root of unity in \(\mathbf{E}\) ; if \(\zeta\) is such a root, then every \(n\) -th root of unity is a power of \(\zeta\) , hence \(\mathbf{E} = \mathbf{K}(\zeta)\) .

For the rest of this No.~we shall choose a separable closure \(K_{s}\) of K. For every \(n \geqslant 1\) prime to p the group \(\mu_{n}(K_{s})\) is cyclic of order n and the field

\[
\mathrm{R} _ {n} (\mathrm{K}) = \mathrm{K} \left(\mu_ {n} \left(\mathrm{K} _ {s}\right)\right)
\]

is a cyclotomic extension of level \(n\) of \(K\) . We can consider \(\mu_n(K_s)\) as a free module of rank 1 over the ring \(Z/nZ\) , and every element \(\sigma\) of \(\text{Gal}(K_s/K)\) induces an automorphism of \(\mu_n(K_s)\) ; there exists therefore a homomorphism \(\chi_n: \text{Gal}(K_s/K) \to (Z/nZ)^*\) characterized by the formula \(u(\zeta) = \zeta^j\) for every \(n\) -th root of unity \(\zeta\) in \(K_s\) , every \(u\) in \(\text{Gal}(K_s/K)\) and every integer \(j\) in the residue class \(\chi_n(u) \mod n\) . Since \(R_n(K) = K(\mu_n(K_s))\) , the kernel of \(\chi_n\) is the subgroup \(\text{Gal}(K_s/R_n(K))\) of \(\text{Gal}(K_s/K)\) ; hence we have \(\chi_n = \varphi_n \circ \psi_n\) where \(\psi_n\) is the restriction homomorphism of \(\text{Gal}(K_s/K)\) over \(\text{Gal}(R_n(K)/K)\) and \(\varphi_n\) is an injective homomorphism of \(\text{Gal}(R_n(K)/K)\) in \((Z/nZ)^*\) . In particular we have the following result:

PROPOSITION 5. --- For every integer \(n \geqslant 1\) prime to \(p\) the extension \(\mathbf{R}_n(\mathbf{K})\) of \(\mathbf{K}\) is abelian of finite degree, its Galois group is isomorphic to a subgroup of \((\mathbf{Z}/n\mathbf{Z})^*\) and its degree divides the order \(\varphi(n)\) of \((\mathbf{Z}/n\mathbf{Z})^*\) .

Let \(\bar{\mathbf{Q}}\) be an algebraic closure of the field \(\mathbf{Q}\) of rational numbers, and let \(n\) be an integer. Then the cyclotomic polynomial \(\Phi_n\) of level \(n\) is defined by

\[
\Phi_ {n} (\mathbf {X}) = \prod_ {\zeta \in S _ {n}} \left(\mathbf {X} - \zeta\right),\tag{5}
\]

where \(S_{n}\) is the set of primitive n-th roots of unity in \(\bar{Q}\) . The polynomial \(\Phi_{n}\) is of degree \(\varphi(n)\) (Prop. 4). It is clear that \(\Phi_{n}(X)\) is invariant under all automorphisms of \(\bar{Q}\) , and so belongs to Q{[}X{]}. Since every element \(\zeta\) of \(S_{n}\) is a root of the polynomial \(X^{n}-1\) , the polynomial \(\Phi_{n}(X)\) divides \(X^{n}-1\) , and the following lemma shows that \(\Phi_{n}(X)\) is a monic polynomial with integer coefficients.

Lemma 2. --- Let \(f, g\) and \(h\) be monic polynomials in \(\mathbf{Q}[\mathbf{X}]\) such that \(f = gh\) . If \(f\) has integer coefficients, the same is true of \(g\) and \(h\) .

Let \(a\) (resp. \(b\) ) be the least of the integers \(\alpha \geqslant 1\) (resp. \(\beta \geqslant 1\) ) such that \(\alpha g\) (resp. \(\beta h\) ) has integer coefficients; we put \(g' = ag\) and \(h' = bh\) and show by reductio ad absurdum that \(a = b = 1\) . If this were not so, there would be a prime divisor \(p\) of \(ab\) . If \(u \in \mathbf{Z}[\mathbf{X}]\) , let us write \(\bar{u}\) for the polynomial with coefficients in the field \(\mathbf{F}_p\) obtained by reduction mod \(p\) of the coefficients of \(u\) . We have \(g'h' = abf\) , whence \(\bar{g}'\bar{h}' = 0\) ; since the ring \(\mathbf{F}_p[\mathbf{X}]\) is an integral domain (IV, p.~9, Prop. 8) we thus have \(\bar{g}' = 0\) or \(\bar{h}' = 0\) . In other words, \(p\) divides all the coefficients of \(g'\) or all those if \(h'\) and this contradicts the hypothesis.

We have the relation

\begin{enumerate}
\def\labelenumi{(\arabic{enumi})}
\setcounter{enumi}{5}
\tightlist
\item
\end{enumerate}

\[
\mathbf {X} ^ {n} - 1 = \prod_ {d \mid n} \Phi_ {d} (\mathbf {X}).
\]

For we have \(\mathbf{X}^n - 1 = \prod_{\zeta \in \mu_n(\mathbf{Q})} (\mathbf{X} - \zeta)\) and the sets \(S_d\) for \(d\) dividing \(n\) form a partition of \(\mu_n(\mathbf{Q})\) .

The formula (6) determines \(\Phi_n(\mathbf{X})\) when we know \(\Phi_d(\mathbf{X})\) for all the divisors \(d < n\) of \(n\) ; since \(\Phi_1(\mathbf{X}) = \mathbf{X} - 1\) , we thus have a recursive procedure for calculating \(\Phi_n\) . For example for prime \(p\) we have

\[
\mathbf {X} ^ {p} - 1 = (\mathbf {X} - 1) \Phi_ {p} (\mathbf {X}),
\]

whence

\begin{enumerate}
\def\labelenumi{(\arabic{enumi})}
\setcounter{enumi}{6}
\tightlist
\item
\end{enumerate}

\[
\Phi_ {p} (\mathbf {X}) = \mathbf {X} ^ {p - 1} + \mathbf {X} ^ {p - 2} + \dots + \mathbf {X} + 1,
\]

and

\[
\Phi_ {p ^ {r + 1}} (\mathbf {X}) = \Phi_ {p} \left(\mathbf {X} ^ {p ^ {r}}\right) \quad \text { for } \quad r \geqslant 0.
\]

Let us list the values of the polynomials \(\Phi_n(\mathbf{X})\) for \(1 \leqslant n \leqslant 12\) :

\[
\Phi_ {1} (\mathbf {X}) = \mathbf {X} - 1
\]

\[
\Phi_ {2} (\mathbf {X}) = \mathbf {X} + 1
\]

\[
\Phi_ {3} (\mathbf {X}) = \mathbf {X} ^ {2} + \mathbf {X} + 1
\]

\[
\Phi_ {4} (\mathbf {X}) = \mathbf {X} ^ {2} + 1
\]

\[
\Phi_ {5} (\mathbf {X}) = \mathbf {X} ^ {4} + \mathbf {X} ^ {3} + \mathbf {X} ^ {2} + \mathbf {X} + 1
\]

\[
\begin{array}{l} \Phi_ {6} (\mathbf {X}) = \mathbf {X} ^ {2} - \mathbf {X} + 1 \\ \Phi_ {7} (\mathbf {X}) = \mathbf {X} ^ {6} + \mathbf {X} ^ {5} + \mathbf {X} ^ {4} + \mathbf {X} ^ {3} + \mathbf {X} ^ {2} + \mathbf {X} + 1 \\ \Phi_ {8} (\mathbf {X}) = \mathbf {X} ^ {4} + 1 \\ \Phi_ {9} (\mathbf {X}) = \mathbf {X} ^ {6} + \mathbf {X} ^ {3} + 1 \\ \Phi_ {1 0} (\mathbf {X}) = \mathbf {X} ^ {4} - \mathbf {X} ^ {3} + \mathbf {X} ^ {2} - \mathbf {X} + 1 \\ \Phi_ {1 1} (\mathbf {X}) = \mathbf {X} ^ {1 0} + \mathbf {X} ^ {9} + \mathbf {X} ^ {8} + \mathbf {X} ^ {7} + \mathbf {X} ^ {6} + \mathbf {X} ^ {5} + \mathbf {X} ^ {4} + \mathbf {X} ^ {3} + \mathbf {X} ^ {2} + \mathbf {X} + 1 \\ \Phi_ {1 2} (\mathbf {X}) = \mathbf {X} ^ {4} - \mathbf {X} ^ {2} + 1 \end{array}
\]

The values of \(\Phi_1, \Phi_2, \Phi_3, \Phi_4, \Phi_5, \Phi_7, \Phi_8, \Phi_9\) and \(\Phi_{11}\) follow directly from (7); now we have \(\Phi_1\Phi_2\Phi_3\Phi_6 = X^6 - 1\) and \(\Phi_1\Phi_2\Phi_3\Phi_4\Phi_6\Phi_{12} = X^{12} - 1\) , whence \(\Phi_4\Phi_{12} = \frac{X^{12} - 1}{X^6 - 1} = X^6 + 1\) and finally \(\Phi_{12} = \frac{X^6 + 1}{X^2 + 1} = X^4 - X^2 + 1\) . The cases \(n = 6\) and \(n = 10\) may be treated similarly.

Remark. --- * For every integer \(n>0\) a function \(\mu(n)\) is defined as follows: if \(n\) is divisible by the square of a prime number, we put \(\mu(n)=0\) , otherwise \(\mu(n)=(-1)^{h}\) if \(n\) is the product of \(h\) distinct prime numbers (« Möbius function »). It may be shown that

\[
\Phi_ {n} (\mathbf {X}) = \prod_ {d \mid n} \left(\mathbf {X} ^ {n / d} - 1\right) ^ {\mu (d)},\tag{8}
\]

or more explicitly

\[
\Phi_ {n} (\mathbf {X}) = \prod_ {p _ {1} <   \dots <   p _ {h}} \left(\mathbf {X} ^ {n / p _ {1} \dots p _ {h}} - 1\right) ^ {(- 1) ^ {h}}\tag{9}
\]

where \((p_{1},\ldots,p_{h})\) runs over the set of all strictly increasing sequences of prime divisors of n (cf.~Lie, II, p.~207, Ex. 1). *

